\phantomsection\addcontentsline{toc}{chapter}{Abstract} % Don't remove me!
\chapter*{\centering{\Large Abstract}}

Bond between a deformed bar and the surrounding concrete governs anchorage
design, yet the development length provisions of AS~3600:2018, Eurocode~2 and
\textit{fib} Model Code 2010 are empirical expressions fitted to test
databases. They return a length but give no access to the stress field that
produces it, and they require recalibration whenever a reinforcement grade
falls outside the range from which they were derived, most recently for the
SENSE600\textsuperscript{®} bars introduced in Australia by InfraBuild. An
analytical treatment carries no such dependence on bar grade.

This thesis develops a closed-form analytical model of a deformed bar bonded
in concrete. The specimen is modelled as three concentric elastic cylinders,
bar, rib layer and concrete; an assumed stress field satisfying equilibrium
and interface continuity reduces the problem to one unknown function, the
axial stress in the bar. Making the strain energy stationary yields a
fourth-order differential equation solved in closed form, the radial-stress
parameter following from a second stationarity condition. No numerical
integration, discretisation or iteration is used. Failure is assessed by two
criteria: yielding of the bar, and cracking of the concrete under a quadratic
interaction of hoop and shear stress averaged over a semicircular region after
Taylor's Theory of Critical Distances. Bond stiffness and slip follow from the
same energy.

The formulation is linear elastic and represents no bond degradation, so it
applies to short embedment lengths, where the stress distribution remains
close to elastic up to failure. It is assessed against pull-out tests rather
than anchorage at full development length. The implementation is verified
against two published analytical solutions, reproducing both without parameter
adjustment, and predictions are compared with pull-out tests from the
literature.

Extending the model to the long embedments that govern development length, and
hence to the anchorage of SENSE600\textsuperscript{®} reinforcement, requires a
softening bond-slip law and an understanding of material non-linearity, and is the principal
direction for further work.